%\input ../util/bookmacros.tex
%\input ../util/docparams.tex

\chapter{Summation Calculus}

\section{Telescoping Sums}
In the last chapter we had a sequence of distance data, $\{x_n\}_{n=0}^{N}$
(this time starting the count at $0$, the morning reading).
If $x_0$ is the odometer reading at the beginning of the day and 
$x_N$ is the reading at the end of the day, then the distance traveled 
that day is $x_N - x_0$. Notice also that the distance traveled over 
the $i^{\rm th}$ interval (our time intervals were every $1/10^{\rm th}$ of a second) is
$x_{i+1} - x_i$. The distance traveled for the whole day should be just 
the sum of all the distances traveled over each small interval. In 
mathematical terms, this would be written:
$$
\obrace{(x_1 - x_0)}{first interval}  + 
\obrace{(x_2 - x_1)}{second interval} + 
\obrace{(x_3 - x_2)}{third interval} + \cdots + 
\obrace{(x_N - x_{N-1})}{Nth interval} = \obrace{(x_N - x_0)}{total distance}
$$
It seems like we have taken something obvious and produced a deep formula.
If you look at the left hand side of the formula you will see that after 
adding the first and second distances the $x_1$'s cancel. If we 
then add the third distance the $x_2$'s cancel. And so on, all the intermediate 
$x$'s cancel except the first and the last. It isn't that mysterious after all. 
We say that such a sum is a 
\index{telescoping sum}; the sum acts like a collapsible telescope.

\section{Fundamental Theorem of Difference Calculus}
Previously, we introduced a notation to describe sequences and their differences; now, since we are talking 
about sums, it is convenient to use a notation for \index{summation}. The symbols
$\sum_{n=1}^{N} x_n$, will mean the sum: $x_1 + x_2 + \cdots + x_N$. 
In our new notation the above \index{telescoping sum} becomes:
$$
\obrace{\sum_{n=0}^{N-1}}{Discrete Sum} \obrace{\Delta x_n}{Telescoping Difference} = 
\obrace{x_{N} - x_0}{Difference at Ends}
$$
{\bf Note:\/} The Greek letter $\sum$ (sigma) corresponds to the letter S -- for sum.

There is nothing special about starting at $0$; the sum telescopes from any
starting index. For any integers $a \le b$ we have
$\sum_{n=a}^{b} \Delta x_n = x_{b+1} - x_a$, and we will use this general
form freely in what follows.

Although the
above formula seemed rather obvious from the way we derived it, we will refer 
to this formula as {\it The Fundamental Theorem of Difference Calculus\/}, or FTDC.
The reason why it is important is that it relates something that is 
computationally hard to do, sum $N$ things, with something that is 
computationally easy, subtract {\it two\/} things. 
This is very useful 
if we have a sum of a large number of things and then recognize that we 
are \index{summing} the \index{difference}s of a sequence. If we make this 
connection, the computation becomes easy. Suppose for instance, that someone 
asks you to sum a sequence of the form: 
$\sum_{n=0}^{N-1} v_n$. If you can find another sequence, $\{x_n\}_{n=0}^{N}$
such that $\Delta x_n = v_n$, then you can sum this series easily by the 
\index{Fundamental Theorem of Difference Calculus} since 
$\sum_{n=0}^{N-1} v_n = \sum_{n=0}^{N-1} \Delta x_n  = x_{N} - x_0$. The 
tricky part, of course, is finding the sequence $\{x_n\}_{n=0}^N$ 
whose difference is $\{v_n\}_{n=0}^{N-1}$. 

In later chapters we will introduce a notion of ``infinitesimal'' differences 
and ``continuous'' sums which will have a similar relationship. 
The corresponding theorem is the ``Fundamental Theorem of Calculus''. 
It is the cornerstone of calculus and it is a generalization of the
simple observation above.

\section{Linearity of Summation and Differencing}
We write down a basic property of both the summing and differencing operations, $\sum$, and $\Delta$.
One of the most {\it important properties\/} of the two operations is that they both 
are {\it \index{Linear Operations}\/} in that they distribute through 
both constant multipliers and sums.
\beginEnum
\enum{$\sum_{n=1}^N (a\, x_n + b\, y_n) = a \sum_{n=1}^N x_n + b \sum_{n=1}^N y_n $.}
\enum{$\Delta (a \,x_n + b \,y_n) = a \,\Delta x_n + b\, \Delta y_n $.}
\endEnum
From this it follows that $\sum_{n=1}^N a x_n = a \sum_{n=1}^N x_n$ and $\Delta (a x_n) = a \Delta x_n$. 
This can be seen by taking $y_n = 0$.

This is the {\it exploit linearity\/} move of the preface making its first
formal appearance: break a problem into pieces, handle each piece, and add
the answers. It allows us to more easily compute sums and differences of
sequences in terms of their component parts. We will need this property almost
immediately, to split differences like $\Delta (n^2 - n)$ into pieces we
already know.%
\footnote{\kern 0.5pt \raise 0.5ex \hbox{\dag}}
{The theme of linearity will continue in the next chapters when we talk about the continuous analogs
of differencing and summing.}

\section{Applications of the Fundamental Theorem}
As an example of how this might be useful consider the problem of summing 
the integers from $1$ to $N$: $\sum_{n=1}^{N} n$. It would be nice if we could find a simple formula
for this sum in terms of $N$. If we could recognize this sum 
as a sum of a \index{telescoping difference} then we could easily find the answer. That is,
if we could find a sequence $\{x_n\}$ such that $\Delta x_n = n$, then 
$\sum_{n=1}^N n = \sum_{n=1}^N \Delta x_n = x_{N+1} - x_1$. We call $x_n$ an {\it anti-difference\/} of the 
sequence $\{n\}_{n=1}^N$. But, can we find such a sequence $\{x_n\}$?
Well, from the last chapter we know that $\Delta n^2 = 2n+1$. We also know that $\Delta n = 1$.
Therefore,  $\Delta (n^2 - n) = \Delta n^2 - \Delta n = (2n + 1) - 1 = 2 n$.%
\footnote{\kern 0.5pt \raise 0.5ex \hbox{\ddag}}{Using the linearity of $\Delta$.}

Finally, if we divide by 2 we have $\Delta \bigl( (n^2 - n)/2 \bigr) = n$.%
\footnote{\kern 0.5pt \raise 0.5ex \hbox{*}}{Linearity of $\Delta$ again.} 
Therefore, if $x_n = (n^2 - n)/2$ then
$\sum_{n=1}^N n = \sum_{n=1}^N \Delta x_n = x_{N+1} - x_1 = ((N+1)^2 - (N+1))/2 - (1^2 - 1)/2 = N (N+1)/2$.
Summarizing:
$$
\sum_{n=1}^N n = N (N + 1)/2
$$
So the sum of the first 1000 numbers is $(1000 * (1000 + 1))/2 = 500,500$.%
\footnote{\kern 0.5pt \raise 0.5ex \hbox{\S}}{For clarity we will occasionally
write ``$a * b$'' instead of the conventional ``$a\, b$'' to
indicate the multiplication of $a$ and $b$.}
One can extend this technique to finding the sum of squares of the first $N$ integers.
The difference of the sequence
$n^3$ is a polynomial of degree 2. So $n^3$ is an imperfect anti-difference in that its difference
produces a multiple of $n^2$ plus lower-order terms -- a term in $n$ and a constant
term. But, as above, if we add combinations of $n^2$ and $n$ to $n^3$ we may
rid ourselves of these extraneous terms.  
That is, we would try to find an {\it \index{anti-difference}\/} of $n^2$ as a linear combination of $n^3$, $n^2$, and $n$.
Once we have this anti-difference we can easily find the formula
for the sum of $n^2$. Let us flesh this out. First, find a linear combination of $n$, $n^2$, and $n^3$ whose 
difference is $n^2$. That is, find $c_1$, $c_2$, and $c_3$ such that 
$$
\Delta (c_1 n + c_2 n^2 + c_3 n^3) = n^2
$$
Since $\Delta$ is linear it distributes over the sum and through multiplication by the $c$'s. That is,
$$
\Delta ( c_1 n + c_2 n^2 + c_3 n^3) = c_1 \Delta n + c_2 \Delta n^2 + c_3 \Delta n^3  = 
c_1 + c_2 (2\, n + 1) + c_3 (3\, n^2 + 3\, n + 1) 
$$
As a polynomial in $n$ this is
$$
(c_1 + c_2 + c_3) + (2 \, c_2 + 3\, c_3) n + 3\, c_3 n^2 
$$
We need to find $c_1$, $c_2$, and $c_3$ so that this is $n^2$. For this to happen the first two terms
should be 0 and the last coefficient should be 1. This gives us three equations:
$$
\eqalign{
c_1 + c_2 + c_3 &= 0 \cr
2\, c_2 + 3 \, c_3 &= 0 \cr
3\, c_3 & = 1
}
$$
Starting at the last equation and working backwards we have $c_3 = 1/3$, $c_2 = -3\, c_3 / 2 = -1/2$, and
$c_1 = -(c_2 + c_3) = 1/6$. So, an anti-difference is $1/6\, n - 1/2\, n^2 + 1/3\, n^3$. Now we can compute
the sum using the FTDC:
$$
\eqalign{
\sum_{n=1}^N n^2 &= \sum_{n=1}^N \Delta (\fracd{1}{6} n - \fracd{1}{2} n^2 + \fracd{1}{3}n^3) \cr
&= (\fracd{1}{6} (N+1) - \fracd{1}{2}(N+1)^2 + \fracd{1}{3} (N+1)^3) - (\fracd{1}{6} * 1 - 
\fracd{1}{2} * 1^2 + \fracd{1}{3} * 1^3) \cr
&= (2\,(N+1)^3 - 3\,(N+1)^2 + (N+1))/6 - 0 \cr
&= (N+1)\bigl[2\,(N+1)^2 - 3\,(N+1) + 1\bigr] / 6\cr
&= (N+1)\bigl[\,((N+1)-1) (2\, (N+1) - 1)\,\bigr]/6 \cr
&= (N+1)(N)(2\, N+1)/6
}
$$
Therefore, the formula is $\sum_{n=1}^N n^2 = N(N+1)(2\, N + 1)/6$.

\section{An Aside: Falling Powers}
The bookkeeping above -- adding combinations of lower powers to kill off the
straggler terms -- suggests that the ordinary powers, $n^k$, are not quite the
natural sequences for differencing. There is a family of sequences that plays
the differencing game perfectly. Define the {\it \index{falling power}\/}
$n^{\underline{k}}$ (read ``$n$ to the $k$ falling'') to be the product of $k$
factors which start at $n$ and step down by one:
$$
n^{\underline{k}} = \obrace{n \, (n-1) \, (n-2) \cdots (n - k + 1)}{$k$ factors}
$$
For example, $n^{\underline{1}} = n$, $n^{\underline{2}} = n(n-1)$, and
$n^{\underline{3}} = n(n-1)(n-2)$. Now compute a difference:
$$
\Delta n^{\underline{2}} = (n+1)\,n - n\,(n-1) = n \bigl( (n+1) - (n-1) \bigr) = 2\, n = 2\, n^{\underline{1}}
$$
The same cancellation works for every $k$. The products $(n+1)^{\underline{k}}$
and $n^{\underline{k}}$ share the $k-1$ factors $n\,(n-1) \cdots (n-k+2)$;
factoring these out leaves $(n+1) - (n-k+1) = k$. That is,
$$
\Delta n^{\underline{k}} = k \, n^{\underline{k-1}}
$$
with {\it no\/} straggler terms -- compare $\Delta n^2 = 2n + 1$ and
$\Delta n^3 = 3n^2 + 3n + 1$. As a one-line application, an anti-difference of
$n = n^{\underline{1}}$ is $n^{\underline{2}}/2$, so by the FTDC
$$
\sum_{n=1}^N n = {(N+1)^{\underline{2}} \over 2} - {1^{\underline{2}} \over 2}
= {(N+1)\,N \over 2} - {1 * 0 \over 2} = {N(N+1) \over 2}
$$
in agreement with the formula found above. When we meet the rule of
differentiation for power functions, $(x^k)^\prime = k\,x^{k-1}$, in the
chapter on Differential Calculus, it is the falling powers that it mirrors
{\it exactly\/}; on the discrete board, the ordinary powers were the wrong
pieces all along.

\exercise{Using the falling power rule, find an anti-difference of
$n^{\underline{2}} = n(n-1)$ and use it, with the FTDC, to compute
$\sum_{n=1}^N n(n-1)$.}

\section{Summation By Parts}
Consider the sum $\sum_{n=1}^{N} n \, 2^n$. The method of the last section
is of no help here: that method hunted for an anti-difference of $n^2$
among linear combinations of the powers $n$, $n^2$, $n^3$, but no
combination of the sequences we have differenced so far has $n \, 2^n$ as
its difference. The trouble is that $n \, 2^n$ is a {\it product\/} of two
very different sequences. We need a tool built for products.

Powerful applications of the \index{FTDC} result from equations that have
a difference on one side. By summing both sides of the equation, the FTDC
implies that the difference side is easy to compute. This means that there
is a simple formula for the summation of the other side of the equation.
We do have such a formula for products; it is the
Discrete Product Rule from the last chapter. It states that:
$$
\Delta (x_n y_n) = \Delta x_n y_{n+1} + x_n \Delta y_n
$$
Summing both sides of this equation from $0$ to $N$ gives us:
$$
\sum_{n=0}^{N} \Delta (x_n y_n) = 
\sum_{n=0}^{N} \Delta x_n y_{n+1} + \sum_{n=0}^{N} x_n \Delta y_n
$$
By the FTDC the left hand side is just $x_{N+1} y_{N+1} - x_0 y_0$. On the right hand 
side we have two complicated sums; moving one of them to the other side of the equation
gives us a way to compare two seemingly unrelated sums.
$$
\sum_{n=0}^{N} x_n \Delta y_n = \bigl(x_{N+1} y_{N+1} - x_0 y_0\bigr) - 
\sum_{n=0}^{N} \Delta x_n y_{n+1}
$$
This formula is referred to as ``Summation by Parts''. It relates one sum
to another. If one sum is known, we can find the other from this formula.
Before unleashing it on the sum $\sum n \, 2^n$ that opened this section,
we first test the new tool by re-deriving two results we already know --
that is how one learns to trust a new tool.

\section{Applications of Summation By Parts}
Let's try to apply \index{summation by parts} to the previous example of 
summing the first $N$ integers. To use ``Summation by Parts'' we 
have to map our problem onto this formula. We can do this if we let 
$x_n = n$ and $\Delta y_n = 1$. We need to find a $y_n$ such that 
$\Delta y_n = 1$; but, we already know that the sequence $y_n = n$ works.  Using the
Summation by Parts formula gives:
$$
\sum_{n=1}^{N} \obrace{n * 1}{$x_n \Delta y_n$} =
\obrace{(N+1) * (N+1)}{$x_{N+1} y_{N+1}$} - \obrace{1 * 1}{$x_1 y_1$}
- \sum_{n=1}^{N} \obrace{1 * (n+1)}{$\Delta x_n y_{n+1}$}
$$
We're trying to use this formula to relate our complicated sum to
a simpler sum. Instead, we end up with a term that is very similar to what we 
started with -- the last term of the above equation. 
This last term can be broken into a piece that is {\it exactly\/} like 
the left hand side of the equation. Unfortunately, it seems like we are 
going backwards. We have
$$
\sum_{n=1}^{N} \obrace{n * 1}{$x_n \Delta y_n$} =
\obrace{(N+1) * (N+1)}{$x_{N+1} y_{N+1}$} - \obrace{1 * 1}{$x_1 y_1$}
- \sum_{n=1}^{N} \obrace{n * 1}{$x_n \Delta y_n$} -
\sum_{n=1}^N 1
$$
But, just as you would when solving an equation in high school algebra 
for ``x'', you collect all expressions of the unknown to one 
side and isolate it; 
everything on the other side is the solution. Our unknown is 
$\sum_{n=1}^{N} n * 1$; that is, $\sum_{n=1}^{N} n$. So, collecting our unknown on 
one side and placing the rest of the expressions on the other yields:
$$
2 \sum_{n=1}^{N} n * 1 = (N+1) * (N+1) - 1 * 1 - \sum_{n=1}^N 1 = 
N^2 + 2N + 1 - 1 - N = N^2 + N
$$
Solving for the sum we have:
$$
\sum_{n=1}^{N} n = N(N+1)/2
$$

\example{Find a formula for the sum: $\sum_{n=1}^{N} n^2$.
We can use the \index{summation by parts} formula again by letting $x_n = n^2$
and finding a $y_n$ such that $\Delta y_n = 1$. $y_n = n$ works as before. 
Therefore,
$$
\sum_{n=1}^{N} n^2 * 1 = (N+1)^2 (N+1) - 1^2 * 1 - 
\sum_{n=1}^{N} \Delta n^2 (n+1)
$$
This simplifies to:
$$
\sum_{n=1}^{N} n^2 * 1 = (N+1)^3 - 1 - 
\sum_{n=1}^{N} (2n+1) (n+1)
$$
Or,
$$
\sum_{n=1}^{N} n^2 * 1 = (N+1)^3 - 1 - 
\biggl( \sum_{n=1}^{N} 2n^2 + \sum_{n=1}^N 3n + \sum_{n=1}^N 1 \biggr)
$$
We have now, as before, an equation for the unknown value $\sum_{n=1}^N n^2$.
Solving for this value gives:
$$
3 \sum_{n=1}^{N} n^2 = (N+1)^3 - 1 - 3 \sum_{n=1}^N n - \sum_{n=1}^N 1
$$
We already know how to sum $\sum_{n=1}^N n$, so we replace it with 
the formula we just worked out.
$$
3 \sum_{n=1}^{N} n^2 = (N+1)^3 - 1 - 3N(N+1)/2 - N
$$
Collecting terms on the right yields
$$
3 \sum_{n=1}^{N} n^2 = N^3 + 3N^2 + 3N + 1 - 1 -\fracd{3}{2}N^2 - 
\fracd{3}{2}N - N
$$
This becomes
$$
3 \sum_{n=1}^{N} n^2 = N\biggl( N^2 + \fracd{3}{2}N + \fracd{1}{2} \biggr) = 
N (N+1)(N+\fracd{1}{2}) = N(N+1)(2N + 1)/2
$$
Finally, we have
$$
\sum_{n=1}^{N} n^2 = N (N+1)(2N+1)/6
$$}

\exercise{Find a formula for the sum: $\sum_{n=1}^{N} n^3$.}

Now for the payoff -- the sum that opened this section, the one the
linear-combination method could not touch.
\example{Find a formula for the sum: $\sum_{n=1}^{N} n \, 2^n$.
Recall from the first chapter that $\Delta 2^n = 2^{n+1} - 2^n = 2^n$; the
sequence $2^n$ is its own difference. So, in the Summation by Parts formula
let $x_n = n$ and $y_n = 2^n$; then $\Delta y_n = 2^n$ and $\Delta x_n = 1$.
The formula gives:
$$
\sum_{n=1}^{N} \obrace{n \, 2^n}{$x_n \Delta y_n$} =
\obrace{(N+1) \, 2^{N+1}}{$x_{N+1} y_{N+1}$} - \obrace{1 * 2}{$x_1 y_1$}
- \sum_{n=1}^{N} \obrace{1 * 2^{n+1}}{$\Delta x_n \, y_{n+1}$}
$$
The remaining sum telescopes for the same reason: $2^{n+1} = \Delta 2^{n+1}$,
so by the FTDC, $\sum_{n=1}^{N} 2^{n+1} = \sum_{n=1}^{N} \Delta 2^{n+1} =
2^{N+2} - 2^2 = 2^{N+2} - 4$. Therefore,
$$
\sum_{n=1}^{N} n \, 2^n = (N+1)\, 2^{N+1} - 2 - \bigl(2^{N+2} - 4\bigr)
= (N+1)\, 2^{N+1} - 2 * 2^{N+1} + 2 = (N-1)\, 2^{N+1} + 2
$$
As a check, for $N = 2$ the formula gives $1 * 8 + 2 = 10$, which agrees
with $1 * 2 + 2 * 4 = 10$.}

\exercise{Find a formula for the sum: $\sum_{n=1}^{N} n^2 \, 2^n$. Hint: in
the Summation by Parts formula let $x_n = n^2$ and $y_n = 2^n$; the sums that
appear can be handled with the formula just derived for
$\sum_{n=1}^{N} n \, 2^n$ and the telescoping of $2^{n+1}$.}


\section{Summary}
The {\it \index{Fundamental Theorem of Difference Calculus}\/} can be stated succinctly:  
The {\it \index{Summation Operation}\/} is the 
inverse of the {\it \index{Difference Operation}\/}. 

This can be seen from the \index{FTDC}:
$\sum_{n=0}^{N-1} \Delta\, x_n = x_N - x_0$, and writing it more crudely as $\sum \Delta\, x = x$. What makes 
this result useful is if for a given sum, $\sum_{n=0}^{N-1} v_n$, one can recognize $v_n$ as the difference of 
another sequence $x_n$, then this sum is $\sum_{n=0}^{N-1} \Delta\, x_n$ and by the FTDC this is just $x_N - x_0$.

The process of recognition is helped greatly if there is a well developed stock of difference formulas. Then
when confronted with summing a sequence, we may be able to recognize, or more easily deduce, a sequence which
is an {\it \index{anti-difference}\/} and then easily compute the sum. Actually, finding differencing and 
anti-differencing formulas for sequences is much like multiplying and dividing numbers. 
Differencing is easier -- like multiplying, while finding anti-differences is harder -- like division. To 
do long division it helps to be fluent with multiplication, because trial and error multiplications are
necessary. The same is true for anti-differencing, it helps first to be fluent with difference formulas for
a number of sequences. To this end, one would
\beginEnum
\enum{Develop difference formulas for a base collection of sequences.}
\enum{Develop difference formulas of sequences formed in various ways from the base; e.g.,
products and {\it \index{Linear Combinations}\/} of sequences.}
\endEnum

This program has been partially carried out in the last chapter. There, we had difference 
formulas for a few sample sequences:
$x_n = c$, $x_n = n$, $x_n = n^2$, $x_n = n^3$. We also had a way of finding 
formulas for sequences formed by linear combinations and 
products of other sequences. 
This helped us above as we tried to find summation formulas for $x_n = n$ and $x_n = n^2$.
We will see in the next two chapters this strategy carried out in more detail. It turns out that
this process often yields simpler formulas when dealing with the continuous analogs of
\index{sequences} (discrete functions) -- \index{continuous function}s.

In fact, the next chapters replay this chapter's game on a bigger board: each
of our discrete tools has a continuous counterpart.
\bigskip
\centerline{\vbox{\halign{\quad\hfil#\hfil\quad&\quad\hfil#\hfil\quad\cr
{\bf Discrete}&{\bf Continuous (coming soon)}\cr
\noalign{\smallskip\hrule\smallskip}
$\Delta$ (difference)&$d/dx$ (derivative)\cr
$\sum$ (summation)&$\int$ (integral)\cr
FTDC&Fundamental Theorem of Calculus\cr
Discrete Product Rule&Product Rule\cr
Summation by Parts&Integration by Parts\cr
}}}
\bigskip


